\scanpage{281}

\unsure{I presume} that in the general case \unsure{where we do} not assume $\Pi_1(\mathcal A^0)=0$, we must replace the condition $\Pi_1(\mathcal A^{0u})=0$ by
\[
\Pi_1(\mathcal A^{0u})\longrightarrow\Pi_1(\mathcal A^0)
\]
an isomorphism.

\margin{Too brutal! The condition in question saying $\Pi_1(\mathcal A^{0\Gamma})\xrightarrow{\sim}\Pi_{\unsure{1}}(\mathcal A^0)^\Gamma$ -- but we must go over all the conditions carefully and modify them ad hoc\ldots{} Cf.\ldots{} \unsure{intercalary|inserted} sheet}

(The nicest thing would be for it to be a \unsure{homotopism} -- it is \unsure{this} probably that would express that we have an $\infty$-\unsure{equivalence}\ldots)

I feel like proving \inserted{first} $1^\circ$, $2^\circ$, $3^\circ$ in the case $\mathcal A=\mathcal A_{g,\nu}$, \unsure{restricting} myself if need be to the \inserted{anabelian} case (probably the hardest in fact! \unsure{\}} but \deleted{more} \inserted{less} conceptually tangled), and of course to the case where $\Gamma$ is finite. The first task will of course be \unsure{that|then} of determining $\mathcal A^{0u}$ and its topological structure, to try to prove that $\mathcal A^{0u}$ is contractible in this case

\scanpage{282}
\mspage{142}

\uuline{24} \inserted{Attempt at} Determination \unsure{of|of the} $\mathcal A^{0\Gamma}$. Link with the relation $\Pi_{g,(\nu,\nu+n-1)}^\Gamma=(1)$, and work program\ldots

Let us consider an operation (faithful \unsure{if we insist}) of the finite group $\Gamma$ on the surface $\unsure{U}\subset X=\widehat U$ of type $(g,\nu)$, \inserted{$U=X-S$,} let $\mathcal A=\operatorname{Aut}(\unsure{U})$, so $\Gamma\subset\mathcal A$ and $\mathcal A^{u}=\mathcal A^\Gamma$ is the centralizer of $\Gamma$. \inserted{Let us first assume $\Gamma=\Gamma^+$.} Obviously $f\in\mathcal A^\Gamma$ implies that $f$ \deleted{\unsure{sends} $\unsure{U}$ into itself} \inserted{induces an \unsure{automorphism} $g$ of $X/\Gamma\unsure{\equiv|\simeq}X'$} \deleted{\unsure{being}} $X$ is a Galois covering with group $\Gamma$.

Let $S'$ be the image of $S$ in $X'$ (so $S=X\mid S'$), and let $S^!$ be the set of ramification points of $\Gamma$ in $U$, $S^{\prime!}$ its image in $U'=X'\backslash S'$ (N.B. $U=X\mid U'$). Thus $g$ is an automorphism of $(X',S')$, taking $S^{\prime!}$ into itself, and \deleted{\unsure{such that}} thus inducing an automorphism of $V'=U'\backslash S^{\prime!}=X'-(S'\cup S^{\prime!})$, \deleted{\unsure{over which}} \inserted{The inverse} \deleted{\illegible{2}} \inserted{image} $V=X\mid V'$ is étale over $V'$, it is even a $\Gamma_{V'}$-torsor over $V'$ and the inverse image of the latter under $g\mid V'$ is (via $f$) isomorphic to $V$. The datum of $f$ is equivalent to the datum of an automorphism of the \deleted{\illegible{1}} $\Gamma$-covering $V$ of $V'$, inducing an

\scanpage{283}

isomorphism of $V'$ which, extended to $X'$, sends $S'$ into itself. This gives an interpretation of $\mathcal A^\Gamma$ in terms of constructions on $X'$.

Let us consider the fundamental group $\Pi'$ of $V'$, \unsure{with the} loop structure indexed by $S'\amalg S^{\prime!}$, so we have a surjective homomorphism
\[
\Pi'\xrightarrow{\varphi}\Gamma
\]
which on \inserted{each} of the subgroups \unsure{with|of} loops of index $s'\in S^{\prime!}$ is not trivial, and an automorphism $g$ of $V$ defines an outer automorphism \inserted{$\unsure{\mathfrak g|\overline g}$} of $\Pi'$ -- the condition to put on it is that its composite with $\varphi$ is conjugate \unsure{to|of} $\varphi$, and that $\unsure{\mathfrak g|\overline g}$ takes $S^{\prime!}$ into itself.

\margin{we choose a point $a\in V$ and its image $a'\in V'$ as base points, to make explicit \illegible{1} $\unsure{\pi|\Pi'}$}

This being so, \deleted{$f$ like \unsure{$g$}} $g$ is defined by $\unsure{\mathfrak g|\overline g}$ up to isotopy \deleted{\unsure{or|and}} (if we leave out degenerate cases \inserted{$(g',\nu')$} \deleted{\illegible{2}}, namely $(0,\nu')$ with $\nu'=0,1,2$ which will have to be examined separately\ldots) and $g$ being chosen, $f$ is defined modulo multiplication by an element of the center of $\Gamma$, but we already know that $\mathcal A^0\cap\Gamma=(1)$, so if we restrict ourselves to examining elements $f\in\mathcal A^{0\Gamma}$ then $f$ will be uniquely determined by $g$

\scanpage{284}
\mspage{143}

We are thus led to the following problem:

Let $V'=X'\setminus(S'\amalg S^{\prime!})$ be an
\deleted{\unsure{orientable}} \inserted{admissible} surface of type $(g',\nu')$
($X'$ compact, $\nu'=\operatorname{card} S'+\operatorname{card} S^{\prime!}$),
equipped with a \unsure{subset} $S^{\prime!}$ of the set of its points at infinity
and a base point $x'$, whence $\Pi'=\Pi_1(\unsure{X'|V'},x')$.
We are given a Galois covering \inserted{\unsure{connected|named} $V$} of $V'$ with
finite group $\Gamma$, pointed over $x'$, hence defined by a surjective homomorphism
\[
\Pi'\xrightarrow{\varphi}\Gamma
\]
and we assume \unsure{$V'$|$V$} ramified at each of the $s'\in S^{\prime!}$,
i.e. (for a choice of the \unsure{group with loops|group of loops} \deleted{\illegible{1}} $L_{s'}$
corresponding to $s'$) that $L_{s'}\to\Gamma$ is not trivial.
We consider the compactified \unsure{set|space} \deleted{\unsure{for|by}} $X$ of $V$,
and the set $S$ resp. $S^!$ of points of $X$ over $S'$ resp. $S^{\prime!}$.
Let $g$ be an automorphism of \unsure{$V$|$V'$}, defining an outer automorphism
$g_\bullet$ of $\Pi'$, we assume that
\[
\unsure{g^*(V'/V)}\overset{\text{\inserted{\unsure{$V$-isomorphism}}}}{\simeq}\unsure{V'}
\quad\text{i.e. that }\varphi\circ g_\bullet\text{ is conjugate to }\varphi
\]
\margin{We may assume that $g$ fixes $x'$, so $g_\bullet$ is a well-defined automorphism of $\Pi'$.}
\unsure{(|\{}which expresses that $g$ comes by passage to the quotient from an automorphism $f$ of $V$,
defined modulo multiplication by an \inserted{\unsure{element of}}
$z\in$ Center $\Gamma$).
Let $U=V\cup S^!=X\setminus S$ (\unsure{ramified covering} over
$U'=V'\cup S^{\prime!}=X'\setminus S'$),
\inserted{\unsure{We want}} \deleted{\illegible{1}} to express that
\deleted{\unsure{$g$ comes by an}} \inserted{among the $f$ which lift}
$g$, there is one (\unsure{necessarily|necessarily} unique!) which is \deleted{in}
\unsure{isotopic} to $1$ in $\underline{\mathrm{Aut}}(U)$ --- i.e. which induces
the trivial outer automorphism of $\Pi_1(U)$. We would like

\scanpage{285}

to prove that these are exactly those which \deleted{\unsure{induce}} \inserted{are}
isotopic to $1$ in $\underline{\mathrm{Aut}}(V')$ --- or again that for such an $f$,
$f$ is necessarily isotopic to $1$ in $\underline{\mathrm{Aut}}(\unsure{V})$
\deleted{\unsure{or $\underline{\mathrm{Aut}}(U)$}}
\inserted{i.e. is in $\underline{\mathrm{Aut}}(\unsure{V})^{0}$
($\subset\underline{\mathrm{Aut}}(\unsure{V|U})^0$)}
\unsure{and} not merely in $\underline{\mathrm{Aut}}(U)^0$.
Let us note \inserted{when $\underline{\mathrm{Aut}}(V')^0$ is $1$-connected} that
$\underline{\mathrm{Aut}}(V')^0$ lifts (by virtue of \unsure{general principles}) to
\deleted{\unsure{$\underline{\mathrm{Aut}}$}}
\deleted{\unsure{$\underline{\mathrm{Aut}}$}}
\[
\underline{\mathrm{Aut}}(V')^0\longrightarrow
\underline{\mathrm{Aut}}(V)^{0\Gamma}\subset
\underline{\mathrm{Aut}}(U)^{0\Gamma}.
\]

Thus the stated condition on $g$ of isotopy to $1$ is certainly
\uline{sufficient}. Let us note moreover that the results already obtained imply that
\[
\underline{\mathrm{Aut}}(U)^{\deleted{\unsure{0}}\Gamma}
\cap\underline{\mathrm{Aut}}(U)^0
=\underline{\mathrm{Aut}}(U)^{0\Gamma}
\]
induces \uline{the identity} on $S^!$ ($=U^!$ \unsure{with} an earlier notation).
Now \deleted{the subgroup} let $B$
\deleted{\unsure{$\subset\underline{\mathrm{Aut}}(U)^0$}}
\inserted{$\subset\underline{\mathrm{Aut}}(U)^0\cap
\underline{\mathrm{Aut}}(\unsure{U|V},S^!)_{\unsure{=|\epsilon}}$}
be the subgroup of points \unsure{which} fix the $s\in S^!$, so that
\[
\underline{\mathrm{Aut}}(U)^0/B\simeq\underline{\mathrm{Mon}}(S^!,U),
\]
and $B^0=\underline{\mathrm{Aut}}(V)^0$ --- a known argument shows us that
\[
B/B^0\simeq\Pi_1(\underline{\mathrm{Mon}}(S^!,U)).
\]
\inserted{Since $\underline{\mathrm{Aut}}(U)^{0\Gamma}\subset B$ i.e.
$\underline{\mathrm{Aut}}(U)^{0\Gamma}=B^\Gamma$, \unsure{we}}
\deleted{We thus have
\diagram{\begin{tikzcd}[ampersand replacement=\&,column sep=1.6em,row sep=1em]
B^{0\Gamma}
 \arrow[r,phantom,"\subset" description]
 \arrow[d,phantom,"{\rotatebox[origin=c]{90}{$\simeq$}}" description]
\& \underline{\mathrm{Aut}}(U)^{0\Gamma}
 \arrow[r,phantom,"=" description]
\& B^\Gamma
\\
\underline{\mathrm{Aut}}(V')^0 \& {} \& {}
\end{tikzcd}}
, \unsure{thus} ($B^{0\Gamma}$ \unsure{is|being} \unsure{connected} and \illegible{2} $B^\Gamma$)}
\deleted{$\underline{\mathrm{Aut}}$}
:
\diagram{\begin{tikzcd}[ampersand replacement=\&]
1 \arrow[r] \& \unsure{B^{0\Gamma^0}} \arrow[r]
  \arrow[d, phantom, "\rotatebox{90}{$\simeq$}" description]
\& B^\Gamma \arrow[r] \arrow[d, phantom, "\rotatebox{90}{$\simeq$}" description]
\& (B/B^0)^\Gamma \arrow[r]
\& H^1(\Gamma,B^0) \arrow[d, equal, "?"] \\
\& \underline{\mathrm{Aut}}(V)^0 \& \underline{\mathrm{Aut}}(U)^{0\Gamma} \&\& 1
\end{tikzcd}}
\margin{\unsure{thus deduce from}
\diagram{\begin{tikzcd}[ampersand replacement=\&]
1\arrow[r]\&B^0\arrow[r]\&B\arrow[r]\&B/B^0\arrow[r]\&1
\end{tikzcd}}}

\scanpage{286}
\mspage{144}

Thus we find that the identity component of
$\underline{\mathrm{Aut}}(U)^{0\Gamma}$ is $\simeq\underline{\mathrm{Aut}}(V)^0$
(\unsure{thus} is $\infty$-connected if $V$ \unsure{anabelian}), and its $\Pi_0$ is
\[
\subset(B/B^0)^\Gamma\simeq\Pi_1(\mathrm{Mon}(S^!,U))^\Gamma.
\]
We would like to prove that the subgroup \unsure{of|of the} \inserted{this} $\Pi_1$ of
invariants under $\Gamma$ is reduced to $1$,
\margin{\unsure{?}}
\deleted{\uline{and better that simultaneously}}
\[
\uline{\unsure{H^0(\Gamma,\Pi_1(\mathrm{Mon}(S^!,U)))}}
\deleted{\uline{=H^1(\Gamma,\Pi_1(\quad))=1}}
\]

We would like this to be true even independently of realizability hypotheses,
when we are given a homomorphism from $\Gamma$ into the Teichmüller group of a $V$,
and we make \unsure{some} suitable hypotheses....

\uline{Conjecture} Let $\Pi$ be an outer group with loops of type $(g,\nu)$,
$I$ the set \unsure{of indices|of index} of the \inserted{classes of} loops,
$\Gamma$ a finite group operating on $\Pi$ by outer automorphisms.
Then there exists an outer group with loops $\Pi'$, \unsure{with set|with an} $I'$ of loop
classes, an outer \unsure{operation|action} of $\Gamma$ on $\Pi'$, a subset
$I^{\prime!}$ of $I'$ stable \unsure{under|for} $\Gamma$, an \inserted{outer} homomorphism
for ``plugging the holes of $I'$''
\[
\Pi'\longrightarrow\Pi
\]
compatible with the action of $\Gamma$, such that

\scanpage{287}

$1^\circ)$ The stabilizer in $\Gamma$ of $\forall$ element of $I^{\prime!}$ is not reduced to $1$

$2^\circ)$ The extension of $\Gamma$ by $\Pi'$ deduced from the outer action of $\Gamma$ has no elements of finite order $\ne 1$ (i.e. for no \unsure{elemen\textsuperscript{t}} of $\Gamma$ $\ne 1$, is the outer action \inserted{on $\Pi'$} realized \unsure{as|by} an automorphism of $\Pi'$ of finite order)

Moreover, the outer $\Gamma$-group with loops $\Pi'$, equipped with the morphism $\Pi'\longrightarrow\Pi$, is determined up to unique isomorphism.

\deleted{\uline{\unsure{Comment.}}} \uline{Comment.} The existence of $\Pi'$, \deleted{and} $\Pi'\longrightarrow\Pi$ is obvious in the ``realizable'' case. Uniqueness -- up to unique isomorphism, \inserted{even} in the realizable case is not obvious, nor even proved. The kernel of \deleted{\illegible{1}}
\diagram{\begin{tikzcd}[ampersand replacement=\&]
\mathrm{Autext}_{\mathrm{lac}}(\Pi',\Pi) \arrow[r] \arrow[d,equal] \& \mathrm{Autext}(\Pi) \\
\mathrm{Autext}_{\mathrm{lac}}(\Pi',I^{\prime!}) \&
\end{tikzcd}}
is precisely the group $\Pi_1$ from earlier
\margin{rather, an extension of \unsure{$\mathfrak{S}_{I'}$} by \unsure{this} $\Pi_1$}
and to show that the functor from the \deleted{\illegible{1}} \inserted{\unsure{pairs} $(\Pi',I^{\prime!})$ of an} \unsure{outer|category} $\Gamma$-group $\Pi'$, \deleted{\illegible{1}} and of a subset

\scanpage{288}
\mspage{145}

$I^{\prime!}$ of $I(\Pi')$ stable under $\Gamma$, such that the corresponding extension of $\Gamma$ by $\Pi'$ is ``torsion-free'' and that the stabilizer of $\forall\ i'\in I^{\prime!}$ in $\Gamma$ is nontrivial, to outer $\Gamma$-groups, is (\unsure{not} only essentially surjective, but) fully faithful, is already problematic. Faithfulness means precisely that $\Pi_1^\Gamma=1$, full faithfulness is \unsure{stronger|stronger} than the fact that \unsure{trivial}
\[
\mathrm{Autext}(\Pi',\Pi)^\Gamma\longrightarrow\mathrm{Autext}(\Pi)^\Gamma
\]
is surjective (in the realizable case, this \inserted{surjectivity} would be a consequence of the fact that $\forall$ \unsure{class} of homeomorphisms commuting with $\Gamma$ contains a homeomorphism commuting with $\Gamma$)
-- we must add to this that \unsure{any|any} other lifting of $\Gamma$ as an outer \unsure{operation|group} of $\Pi'$ [i.e. a $\Gamma\longrightarrow\mathrm{Autext}_{\mathrm{lac}}(\Pi',\Pi)$\inserted{\diagram{\begin{tikzcd}[ampersand replacement=\&,row sep=1.2em]
\mathrm{Autext}_{\mathrm{lac}}(\Pi',\Pi)
 \arrow[d,phantom,"\inserted{\unsure{\rotatebox[origin=c]{90}{$\simeq$}}}" description]
\\
\inserted{\mathrm{Autext}_{\mathrm{lac}}(\Pi',I)}
\end{tikzcd}}}], \deleted{\unsure{such that}} \inserted{having the} \unsure{$=$|same} \uline{action of $\Gamma$} \uline{on} $I^{\prime!}$, is conjugate \unsure{by|to} the \illegible{1}

\scanpage{289}
\unsure{previous one|previous one} by an element of the \unsure{kernel} $\Pi_1$ of
\[
\mathrm{Autext}_{\unsure{\mathrm{lac}|\mathrm{loc}}}(\Pi',I^{\prime!})
\longrightarrow
\mathrm{Autext}_{\unsure{\mathrm{lac}|\mathrm{loc}}}(\Pi)
\times \unsure{\mathfrak S}_{I^{\prime!}}\quad\ldots
\]

We are going to \unsure{restate} the preceding conjecture in a slightly more general form.
We are going to introduce a notation. If $\Gamma$ is a finite group operating \unsure{by outer automorphisms}
\inserted{in an orientation-preserving way} on a group with loops $\Pi$ \inserted{\unsure{anabelian}},
giving rise to an extension $\unsure{\mathcal E|E}$ of $\Pi$ by $\Gamma$,
we \unsure{will} \unsure{denote} by $\Phi=\Phi(\Pi,\Gamma)$ (``fixed \unsure{points}'')
\unsure{the set} of $\Pi$-conjugacy classes of maximal partial sections
\inserted{$\ne(1)$} of the extension.
(The intrinsic character of $\Phi(\Pi,\Gamma)$, independently of the particular choices of $\unsure{\mathcal E|E}$, has already been noted).
It is clear that $\Phi$ is a $\Gamma$-set, we also know that it is finite.

Let \unsure{now} $(\Pi',\Gamma)$ be an \inserted{\unsure{outer}} $\Gamma$-group
with loops, \inserted{\unsure{with} $\Gamma$ operating in an orientation-preserving way ($\Gamma=\Gamma^+$)
and \unsure{faithfully} (\unsure{to simplify})} \unsure{let us fix} a subset
$I^{\prime!}$ of $I'=I(\Pi')$, stable under $\Gamma$,
\inserted{such that $\forall i'\in I^{\prime!}$, \unsure{we} have $\Gamma_{i'}\ne(1)$,}
and consider the quotient outer group, defined by plugging $I^{\prime!}$,
namely $\Pi$ \deleted{\unsure{$\Pi'_{I^{\prime!}}$}}, -- it is clear that $\Gamma$ still operates on it.
We want first to define a bijection

\scanpage{290}
\mspage{146}
\[
(*)\qquad\Phi(\Pi,\Gamma)\overset{\sim}{\longleftarrow}
\Phi(\Pi',\Gamma)\unsure{\amalg|\sqcup} I^{\prime!}
\]
\margin{assuming $\Pi,\Gamma$ also
\unsure{anabelian}, to be more sure!}

a) Map $I'\longrightarrow\Phi(\Pi,\Gamma)$

Let $L_{i'}\subset\Pi'$ be of \unsure{class} $i'\in I'$, $Z_{i'}$ its centralizer in the extension $\unsure{\mathcal E|E}'$ of $\Gamma$ by $\Pi'$
so that $L_{i'}$ is an extension of $\Gamma_{i'}$
by $L_{i'}$. The passage to the quotient $\unsure{\mathcal E|E}'\longrightarrow\unsure{\mathcal E|E}$
defines a partial section $\Gamma_{i'}\hookrightarrow\unsure{\mathcal E|E}$,
contained in a unique \unsure{maximal} partial section
\inserted{(probably it is already maximal \unsure{$=$|itself}?\ldots)}.
\margin{Yes, because for \illegible{2} \unsure{$\Gamma$ cyclic},
\illegible{5} \unsure{realizable}\ldots ! \unsure{?}}
We thus find a map $I'\longrightarrow\Phi(\Pi,\Gamma)$,
obviously compatible with the actions of $\Gamma$.

b) Map $\Phi(\Pi',\Gamma)\longrightarrow\Phi(\Pi,\Gamma)$

Every partial section \inserted{$\ne1$} of $\unsure{\mathcal E|E}'$ over $\Gamma$ obviously defines
one of $\unsure{\mathcal E|E}$ over $\Gamma$ by passing to the
quotient, we know as above that a
maximal section gives a maximal section,
by reducing to the case $\Gamma$ cyclic.
\margin{not clear, since \unsure{the hypothesis} that the $\Gamma_{i'}\ne(1)$
\unsure{risks|risk} not being \unsure{satisfied|satisfied}}

c) Bijectivity of (*). \deleted{For surjectivity,} \inserted{It is not}
\deleted{we are \unsure{reduced} to the case}
\inserted{\unsure{proved} (\unsure{except} if the initial situation is}
realizable, or if $I^{\prime!}$ reduced to a single element
\ldots but then $\Gamma$ cyclic and the situation is realizable\ldots)